\documentclass[12pt,a4paper]{article}
\usepackage[utf8]{inputenc}
\usepackage[spanish]{babel}
\usepackage{amsmath,amssymb}
\usepackage{geometry}
\geometry{margin=2.5cm}
\usepackage{booktabs}
\usepackage{array}
\usepackage{graphicx}
\usepackage{hyperref}
\usepackage{listings}
\usepackage{xcolor}
\usepackage{caption}
\usepackage{subcaption}
\usepackage{float}
\usepackage{enumitem}
\usepackage{fancyhdr}
\usepackage{multirow}
\usepackage{pgfplots}
\pgfplotsset{compat=1.18}
\setlength{\headheight}{15pt}

\pagestyle{fancy}
\fancyhf{}
\fancyhead[L]{INFO1195 --- Actividad 4}
\fancyhead[R]{Filtro Sobel CUDA}
\fancyfoot[C]{\thepage}

\lstset{
  basicstyle=\ttfamily\small,
  backgroundcolor=\color{gray!10},
  frame=single,
  breaklines=true,
  columns=fullflexible
}

\title{\textbf{Implementaci\'on y Comparaci\'on de una Pipeline de\\ Procesamiento de Im\'agenes RGB en GPU mediante CUDA}}
\author{
  Christian Ferrer, Sebasti\'an Cisternas, Patricio Vald\'es, \\
  Eliaser Concha y Benjam\'in Borgeaud \\
  \vspace{0.3cm}
  \textit{Actividad 4 --- INFO1195}
}
\date{Julio 2026}

\begin{document}

\maketitle
\thispagestyle{empty}

\begin{abstract}
Este informe presenta la implementaci\'on, an\'alisis y comparaci\'on de una pipeline completa de procesamiento de im\'agenes RGB sobre GPU, contrastando cuatro versiones: CPU secuencial, CUDA C++ cl\'asico, CUDA Tile C++ y cuTile Python. La pipeline aplica conversi\'on a luminancia, filtro gaussiano, operador de Sobel y redimensionamiento bilineal. Se implement\'o adem\'as una variante de convoluci\'on gaussiana separable (dos pases 1D) para comparar contra la versi\'on directa 2D. Los experimentos se ejecutaron en dos arquitecturas: una \textbf{NVIDIA RTX 4090} (servidor universitario) como entorno principal, y una \textbf{NVIDIA RTX 3060} (equipo local) como l\'inea base. Se reportaron speed-ups de hasta \textbf{167.5$\times$} en la RTX 4090 del servidor y \textbf{156.5$\times$} en la RTX 3060 local para im\'agenes grandes con kernel $9 \times 9$, validando que la GPU es la plataforma \'optima para este tipo de procesamiento. La versi\'on separable redujo el tiempo de kernel del blur en un 25\% adicional, alcanzando 44,082 MP/s de throughput en la 4090.
\end{abstract}

\newpage
\tableofcontents
\newpage

%=============================================================================
\section{Introducci\'on, Contexto y Formulaci\'on de la Pipeline}
%=============================================================================

\subsection{Problema y Motivaci\'on}

% El procesamiento digital de im\'agenes aplica la misma operaci\'on sobre millones de p\'ixeles independientes, lo que lo convierte en un dominio ideal para estudiar paralelismo en GPU. Esta actividad implementa una pipeline de cuatro etapas sobre cuatro versiones comparables, siguiendo el modelo de programaci\'on SIMT (Single Instruction, Multiple Threads) de CUDA.

El procesamiento digital de im\'agenes es una tarea fundamental en aplicaciones como visión por computadora, robótica y sistemas de seguridad. Dónde, estas aplicaciones requieren ejecutar operaciones como filtro gaussiano, detección de bordes y redimensionamiento de imágenes de alta resolución, lo que implica un procesamiento de millones de píxeles mediante determinados cálculos. Al ejecutar estas operaciones de manera secuencial por CPU, el tiempo de procesamiento aumenta considerablemente a medida que crece el tamaño de la imagen o la complejidad de las operaciones, limitando su rendimiento en aplicaciones que requieren respuestas rápidas o procesamiento de grandes volúmenes de datos.

Debido a que cada píxel se puede procesar de manera independiente, este tipo de problemas presenta un desafío que se puede resolver mediante el paralelismo de datos, lo que lo convierte en un caso de estudio adecuado para arquitecturas GPU. Sin embargo, existen diferentes versiones o modos para explotar dicho paralelismo, cada uno con sus características en cuanto al manejo de hilos, organización de la memoria, acceso a datos y complejidad de implementación.

La motivación que da lugar a este trabajo radica en estudiar el impacto de estas versiones sobre el rendimiento de una pipeline completa de procesamiento de imágenes en planos RGB. Para ello, se implementan las etapas de conversión a escala de grises o luminancia, filtro gaussiano, operador de Sobel y redimensionamiento bilineal usando las versiones secuencial en CPU, CUDA clásico en C++, CUDA Tile en C++ y una versión parcial del procesamiento usando cuTile en Python. Finalmente, se comparan sus tiempos de ejecución, throughput, speed-up, uso de la memoria y comportamiento experimental, con la finalidad de analizar las ventajas y limitaciones de cada alternativa de implementación sobre GPU.

\subsection{Pipeline de Procesamiento}

% La pipeline recibe una imagen RGB (PNG) y produce im\'agenes de salida por cada etapa intermedia:

La pipeline implementada corresponde a un proceso de 4 etapas aplicadas sobre una imagen en el plano RGB como entrada. Cada etapa recibe como entrada el resultado de la etapa anterior y produce una nueva imagen que sirve de entrada para la siguiente operación. Este flujo separa responsabilidades y facilita la validación de cada algoritmo.

Cabe destacar que la imagen de entrada se almacena en un formato RGB interleaved o intercalado, dónde cada píxel está compuesto por sus componentes rojo, verde y azul, y a partir de esta representación se ejecutan las siguientes etapas:

\begin{enumerate}
  \item \textbf{Conversión RGB a luminancia y escala de grises:} esta etapa transforma la imagen RGB en una imagen de un solo plano o canal de intensidad llamada escala de grises o luminancia. La cual calcula una combinación ponderada de los canales rojo, verde y azul usando los coeficientes del estándar ITU-R BT.601: $Y = 0.299R + 0.587G + 0.114B$. La finalidad de esta etapa es reducir la cantidad de información que se debe procesar en las siguientes etapas y eliminar la influencia del color, ya que tanto el filtro gaussiano como el de Sobel trabajan únicamente sobre la intensidad de la imagen.
  \item \textbf{Filtro gaussiano:} dada una imagen en escala de grises como entrada, se aplica un filtro de desenfoque gaussiano cuyo objetivo busca reducir el ruido y suavizar pequeñas variaciones de intensidad. El filtro se implementa mediante convolución utilizando una configuración de 5x5 o 9x9, dependiendo del diseño experimental.
  \item \textbf{Filtro de Sobel:} esta etapa realiza la detección de bordes, dónde cada píxel se analiza una vecindad de 3x3, calculando dos gradientes. La gradiente $G_{x}$ permite medir los cambios horizontales de intensidades, mientras que $G_{y}$ mide los cambios verticales. Luego, ambos gradientes se combinan para obtener la magnitud del borde, generando como salida una imagen donde los cambios bruscos de intensidad aparecen resaltados, mientras que las regiones homogéneas permanecen oscuras.
  \item \textbf{Redimensionamiento bilineal:} esta etapa redimensiona una imagen de entrada mediante interpolación bilineal. Para cada píxel de salida se calcula la posición correspondiente dentro de la imagen original mediante coordenadas fraccionales, las cuales utilizan los cuatro vecinos más cercanos para estimar un nuevo valor de intensidad mediante una interpolación ponderada. La implementación de esta etapa busca trabajar con escalas de 0.5 y 1.75 para evaluar tanto la reducción como la ampliación de la imagen.
\end{enumerate}

\begin{figure}[H]
\centering
\fbox{\begin{tabular}{c}
Imagen RGB (PNG) \\
$\downarrow$ \\
{[}1{]} Escala de grises (luminancia 0.299/0.587/0.114) \\
$\downarrow$ \\
{[}2{]} Desenfoque gaussiano ($5\times5$ o $9\times9$, directo o separable) \\
$\downarrow$ \\
{[}3{]} Operador de Sobel ($G_x$, $G_y$, magnitud saturada) \\
$\downarrow$ \\
{[}4{]} Redimensionamiento bilineal ($0.5\times$ o $1.75\times$) \\
$\downarrow$ \\
Im\'agenes de salida por etapa $+$ CSV de m\'etricas
\end{tabular}}
\caption{Diagrama de la pipeline de procesamiento}
\end{figure}

\subsection{Versiones Implementadas}

\begin{table}[H]
\centering
\caption{Versiones de la pipeline}
\begin{tabular}{lll}
\toprule
\textbf{\#} & \textbf{Versi\'on} & \textbf{Toolchain} \\
\midrule
1 & CPU secuencial (referencia) & \texttt{nvcc -std=c++17} \\
2 & CUDA C++ cl\'asico (sin Tile) & \texttt{nvcc -std=c++17 -arch=native} \\
3 & CUDA Tile C++ & \texttt{nvcc -std=c++20 -enable-tile} \\
4 & cuTile Python & \texttt{cuda.tile 1.4.0} + \texttt{torch 2.12.1+cu130} \\
\bottomrule
\end{tabular}
\end{table}

%=============================================================================
\section{Fundamento Matem\'atico}
%=============================================================================

\subsection{Filtro Gaussiano}

El kernel gaussiano 2D se define como:

\begin{equation}
G(x, y) = \frac{1}{2\pi\sigma^2} \exp\!\left(-\frac{x^2 + y^2}{2\sigma^2}\right)
\end{equation}

La implementaci\'on directa realiza una convoluci\'on 2D completa:

\begin{equation}
I_{\text{out}}(x,y) = \sum_{k_y=-r}^{r} \sum_{k_x=-r}^{r} I_{\text{in}}(x+k_x,\, y+k_y) \cdot G(k_x, k_y)
\end{equation}

Complejidad: $O(k^2)$ por p\'ixel, donde $k$ es el tama\~no del kernel.

La variante \textbf{separable} aprovecha que el kernel gaussiano es separable radialmente, $G(x,y) = G_{1D}(x) \cdot G_{1D}(y)$, donde:

\begin{equation}
G_{1D}(x) = \frac{1}{\sqrt{2\pi}\,\sigma} \exp\!\left(-\frac{x^2}{2\sigma^2}\right)
\end{equation}

Realizando dos pases 1D consecutivos:

\begin{align}
I_{\text{temp}}(x,y) &= \sum_{k_x=-r}^{r} I_{\text{in}}(x+k_x,\, y) \cdot G_{1D}(k_x) \quad \text{(horizontal)} \\
I_{\text{out}}(x,y)  &= \sum_{k_y=-r}^{r} I_{\text{temp}}(x,\, y+k_y) \cdot G_{1D}(k_y) \quad \text{(vertical)}
\end{align}

Complejidad: $O(2k)$ por p\'ixel. Para $k=5$, reduce de 25 a 10 operaciones por p\'ixel.

En todas las versiones se usa \textbf{clamp} (r\'eplica del p\'ixel del borde) como estrategia de bordes, y el kernel se normaliza para que su suma sea exactamente 1.

Con respecto a la selección de sigma, en esta implementación se calcula como:

\[
\sigma = \frac{k-1}{6},
\]

donde \(k\) corresponde al tamaño del kernel gaussiano. Esta expresión se basa en la propiedad de la distribución normal de que aproximadamente el \(99.7\%\) de sus valores se encuentran dentro del intervalo \(\pm 3\sigma\). Considerando que el radio del kernel es

\[
r = \frac{k-1}{2},
\]

se busca coincidir dicho radio con tres desviaciones estándar, es decir,

\[
3\sigma = r.
\]

Despejando \(\sigma\), se obtiene

\[
\sigma = \frac{r}{3}
       = \frac{(k-1)/2}{3}
       = \frac{k-1}{6}.
\]

Esta decisión de diseño, busca garantizar que el tamaño del kernel y el grado de suavizado sean consistentes entre sí, evitando ciertas configuraciones poco adecuadas que pueden afectar o degradar la calidad del filtro. Por otro lado, al establecer una relación entre el tamaño del kernel y el cálculo adecuado de sigma, los experimentos son más reproducibles y comparables haciendo que el comportamiento del filtro quede determinado únicamente por el tamaño del kernel seleccionado.

El tamaño del kernel gaussiano debe ser un número impar para que exista un píxel central alrededor del cual se define la vecindad de la convolución. Este píxel actúa como referencia para aplicar simétricamente los coeficientes de la máscara gaussiana sobre los píxeles vecinos.

\subsection{Operador de Sobel}

Las m\'ascaras de convoluci\'on $3 \times 3$ para detectar bordes horizontales y verticales son:

\begin{equation}
G_x = \begin{pmatrix}-1 & 0 & +1\\-2 & 0 & +2\\-1 & 0 & +1\end{pmatrix}, \quad
G_y = \begin{pmatrix}-1 & -2 & -1\\0 & 0 & 0\\+1 & +2 & +1\end{pmatrix}
\end{equation}

La magnitud del gradiente se calcula y se satura a 8 bits:

\begin{equation}
G = \min\!\left(\left\lfloor\sqrt{G_x^2 + G_y^2}\right\rfloor,\; 255\right)
\end{equation}

El operador se aplica sobre la imagen en escala de grises, siguiendo la recomendaci\'on del enunciado. Bordes con clamp.

\subsection{Interpolaci\'on Bilineal}

Para un p\'ixel de salida $(o_x, o_y)$, se calcula la coordenada fraccional en la imagen de entrada:

\begin{equation}
s_x = o_x \cdot \frac{W_{\text{in}}}{W_{\text{out}}}, \quad s_y = o_y \cdot \frac{H_{\text{in}}}{H_{\text{out}}}
\end{equation}

Con los cuatro vecinos $(x_0,y_0) = (\lfloor s_x\rfloor, \lfloor s_y\rfloor)$, $x_1=\min(x_0+1, W_{\text{in}}-1)$, $y_1=\min(y_0+1, H_{\text{in}}-1)$ y diferencias fraccionales $d_x = s_x-x_0$, $d_y = s_y-y_0$:

\begin{equation}
V = A(1-d_x)(1-d_y) + B\cdot d_x(1-d_y) + C(1-d_x)d_y + D\cdot d_x d_y
\end{equation}

\subsection{Conversi\'on RGB $\rightarrow$ Luminancia}

\begin{equation}
Y = 0.299\,R + 0.587\,G + 0.114\,B
\end{equation}

Coeficientes del est\'andar ITU-R BT.601. Aplicada p\'ixel a p\'ixel sobre la imagen RGB de entrada.

\subsection{Validaci\'on Num\'erica}

La correctitud se valid\'o comparando p\'ixel a p\'ixel las salidas de CPU contra CUDA sobre \texttt{small/pistola.png} ($k=5$, $s=0.5$):

\begin{table}[H]
\centering
\caption{Error cuadr\'atico medio (MAE) CPU vs CUDA}
\begin{tabular}{lcc}
\toprule
\textbf{Etapa} & \textbf{MAE} & \textbf{Diferencia m\'axima} \\
\midrule
Grayscale & 0.0098 & 3 \\
Blur & 0.0101 & 3 \\
Sobel & 0.0744 & 15 \\
Resize & 0.0743 & 15 \\
\bottomrule
\end{tabular}
\end{table}

Las diferencias provienen de la promoci\'on a \texttt{float} en GPU y del orden de operaciones en la ra\'iz cuadrada. Tolerancias de la r\'ubrica cumplidas con margen (MAE $\leq$ 1.0 para grises, $\leq$ 2.0 para Sobel).

%=============================================================================
\section{Metodolog\'ia y Dise\~no Experimental}
%=============================================================================

\subsection{Im\'agenes de Entrada}

\begin{table}[H]
\centering
\caption{Im\'agenes de prueba}
\begin{tabular}{lccc}
\toprule
\textbf{Instancia} & \textbf{Archivo} & \textbf{Dimensiones} & \textbf{P\'ixeles (M)} \\
\midrule
Small & \texttt{pistola.png} & $512 \times 512$ & 0.26 \\
Medium & \texttt{ak-47.png}  & $2048 \times 2048$ & 4.19 \\
Large & \texttt{m16.png}    & $4096 \times 4096$ & 16.78 \\
No-divisible & \texttt{automovil.png} & $1402 \times 1122$ & 1.57 \\
\bottomrule
\end{tabular}
\end{table}

La instancia \texttt{no-divisible} ($1402 \times 1122$) no es m\'ultiplo de 16, 32 ni 64, lo que permite validar el manejo de dimensiones no alineadas a tiles.

\subsection{Par\'ametros Experimentales}

\begin{table}[H]
\centering
\caption{Configuraci\'on experimental}
\begin{tabular}{ll}
\toprule
\textbf{Par\'ametro} & \textbf{Valores} \\
\midrule
Tama\~nos de kernel gaussiano & $5 \times 5$, $9 \times 9$ \\
Factores de resize & $0.5\times$, $1.75\times$ \\
Repeticiones por configuraci\'on & 10 \\
Variantes de blur (solo CUDA) & Directa 2D, Separable 2-pases \\
Tama\~no de bloque Tile C++ & $16 \times 16$ (fijo) \\
Tama\~no de bloque CUDA cl\'asico & Din\'amico (Occupancy API) \\
\bottomrule
\end{tabular}
\end{table}

\subsection{Hardware}

Para evaluar la escalabilidad de la pipeline, las pruebas se ejecutaron en dos arquitecturas: el servidor oficial de la universidad como entorno principal, y un equipo local como l\'inea base de comparaci\'on.

\begin{table}[H]
\centering
\caption{Especificaciones del hardware}
\begin{tabular}{lll}
\toprule
\textbf{Componente} & \textbf{Servidor Universidad} & \textbf{Equipo Local} \\
\midrule
GPU & NVIDIA RTX 4090 & NVIDIA RTX 3060 \\
VRAM & 24 GB GDDR6X & 12 GB GDDR6 \\
CUDA Cores & 16,384 & 3,584 \\
Compute Capability & 8.9 & 8.6 \\
CPU & Intel i9-14900KF & Ryzen 5 5600X \\
RAM & 125 GB DDR5 & 32 GB DDR4 \\
CUDA Toolkit & 13.3 & 13.3 \\
Compiler & nvcc + g++ & nvcc + MSVC 2022 \\
\bottomrule
\end{tabular}
\end{table}

\subsection{Herramientas de Medici\'on}

\begin{itemize}
  \item \textbf{CUDA Events:} 8 \texttt{cudaEvent\_t} por etapa (H$\rightarrow$D, kernel, D$\rightarrow$H, total) para separar transferencia de c\'omputo en CUDA cl\'asico y Tile C++.
  \item \textbf{CPU:} \texttt{std::chrono::high\_resolution\_clock} por etapa.
  \item \textbf{cuTile Python:} \texttt{torch.cuda.Event} con \texttt{enable\_timing=True}.
  \item \textbf{Profiling:} Nsight Systems 2026.1.3 para l\'inea de tiempo completa de transferencias y kernels.
\end{itemize}

\subsection{Procedimiento}

\begin{enumerate}
  \item Compilar los 3 binarios C++: \texttt{make} (Linux) o \texttt{build.bat} (Windows).
  \item Ejecutar \texttt{run\_all.sh} (Linux) o \texttt{run\_all.bat} (Windows): 4 instancias $\times$ 2 kernels $\times$ 2 escalas $\times$ 10 repeticiones para cada una de las 5 configuraciones (CPU, CUDA directo, CUDA separable, Tile, cuTile Python).
  \item El CSV unificado se almacena en \texttt{results/resultados.csv} (800 filas de datos).
  \item Profiling con Nsight Systems para an\'alisis de cuellos de botella.
\end{enumerate}

%=============================================================================
\section{Descripci\'on T\'ecnica de las Versiones}
%=============================================================================

\subsection{CPU Secuencial (\texttt{image.cpp}, \texttt{main.cpp})}

Implementaci\'on en C++17 sin SIMD expl\'icto. El kernel gaussiano se genera en formato 2D (\texttt{float**}) y se aplica con convoluci\'on directa. La medici\'on usa \texttt{std::chrono::}\-\texttt{high\_resolution\_clock} por etapa. Diez repeticiones con iteraci\'on de warmup. Las columnas de transferencia H$\leftrightarrow$D quedan en 0 (no hay GPU).

\subsection{CUDA C++ Cl\'asico (\texttt{cuda\_kernels.cu}, \texttt{main\_cuda.cpp})}

Cuatro kernels \texttt{\_\_global\_\_} con grilla 2D, borde clamp y acceso coalescente a memoria global. El tama\~no de bloque se elige din\'amicamente con la API de Occupancy de CUDA (\texttt{cudaOccupancy\-MaxPotential\-BlockSize}). Cada etapa registra 4 mediciones con \texttt{cudaEvent\_t}: H$\rightarrow$D, kernel, D$\rightarrow$H y total. Se implementan dos variantes del blur:

\begin{itemize}
  \item \textbf{Directa 2D:} un kernel que aplica el kernel completo en una sola pasada ($O(k^2)$).
  \item \textbf{Separable 2-pases:} dos kernels consecutivos (horizontal + vertical) con un buffer intermedio en GPU. El kernel 1D se calcula en CPU con \texttt{createSeparableGaussKernel()} y se copia a GPU. Complejidad $O(2k)$. Se activa con el flag \texttt{--separable}.
\end{itemize}

\subsection{CUDA Tile C++ (\texttt{tile\_kernels.cu}, \texttt{main\_tile.cpp})}

Compilado con \texttt{-enable-tile -std=c++20}. Contiene un kernel \texttt{\_\_tile\_global\_\_} real (\texttt{tileIdentityKernel}, $16\times16$, load/store masked) que demuestra que el toolchain Tile est\'a habilitado. Los 4 kernels principales de la pipeline corren como SIMT en el mismo \texttt{.cu} con \texttt{BLOCK\_SIZE = 16} fijo. Esta decisi\'on se documenta en el informe porque la API \texttt{cuda\_tile.h} de CUDA 13.3 no soporta stencils con acceso a vecinos de forma pr\'actica; mezclar Tile + SIMT en la misma unidad de traducci\'on est\'a oficialmente soportado.

\subsection{cuTile Python (\texttt{cutile\_pipeline.py})}

La etapa de RGB $\rightarrow$ grayscale se implementa con \texttt{@ct.kernel} usando tiles 2D separados para los canales R, G y B. Los tiles se cargan con \texttt{PaddingMode.ZERO} en los bordes (a diferencia del clamp usado en las dem\'as versiones), pero esto no afecta la salida porque los p\'ixeles de padding no se almacenan. El resto de etapas (blur, Sobel, resize) usan operaciones gen\'ericas de PyTorch sobre GPU: \texttt{F.conv2d} con kernels calculados manualmente desde la f\'ormula matem\'atica, y \texttt{F.interpolate} con modo bilineal. No se utilizan filtros pre-implementados de PyTorch u OpenCV. El padding en blur y Sobel se configura con \texttt{mode='replicate'} (equivalente a clamp).

\subsection{Layout de Memoria e Indexaci\'on}

\begin{itemize}
  \item \textbf{CPU:} buffer plano \texttt{unsigned char*} de 1 canal (grayscale) o 3--4 canales (RGB/RGBA). \'Indice: $y \cdot W + x$ para gray, $(y \cdot W + x) \cdot C + c$ para RGB.
  \item \textbf{GPU:} mismo layout interleaved. Los kernels leen canales en el offset $(y \cdot W + x) \cdot C + c$.
  \item \textbf{cuTile:} 3 tensores 2D separados para R, G, B (planar) en GPU, cargados como tiles y combinados con $0.299R + 0.587G + 0.114B$.
\end{itemize}

%=============================================================================
\section{Resultados, Tablas y Gr\'aficos}
%=============================================================================

Para evaluar la solidez de la pipeline y su escalabilidad, las pruebas principales se ejecutaron en el servidor universitario (RTX 4090), utilizando el equipo local (RTX 3060) como l\'inea base comparativa. Los resultados se obtuvieron con 10 repeticiones por configuraci\'on.

\subsection{Resultados Principales en Servidor (RTX 4090, $k=5$, $s=0.5$)}

\begin{table}[H]
\centering
\caption{Tiempos y speed-up en el servidor universitario (RTX 4090 + i9-14900KF)}
\label{tab:server_main}
\begin{tabular}{lcccccc}
\toprule
\textbf{Instancia} & \textbf{P\'ixeles} & \textbf{CPU (ms)} & \textbf{CUDA (ms)} & \textbf{Tile (ms)} & \textbf{Sp CUDA} & \textbf{Sp Tile} \\
\midrule
Small    & 0.26M & 19.85 $\pm$ 0.2  & 0.63 $\pm$ 0.03 & 0.61 $\pm$ 0.02 & 31.7$\times$ & 32.5$\times$ \\
Medium   & 4.19M & 317.32 $\pm$ 1.6 & 4.24 $\pm$ 0.28 & 4.19 $\pm$ 0.25 & 74.9$\times$ & 75.8$\times$ \\
Large    & 16.8M & 1333.27 $\pm$ 2.5& 14.55 $\pm$ 0.70& 14.50 $\pm$ 0.84& 91.7$\times$ & 91.9$\times$ \\
No-div   & 1.57M & 119.60 $\pm$ 0.4 & 2.08 $\pm$ 0.08 & 2.08 $\pm$ 0.09 & 57.4$\times$ & 57.4$\times$ \\
\bottomrule
\end{tabular}
\end{table}

El speed-up escala con el tama\~no de imagen: de 31.7$\times$ en $512^2$ a 91.7$\times$ en $4096^2$. La GPU aprovecha m\'as paralelismo a mayor cantidad de p\'ixeles disponibles. CUDA Tile es marginalmente m\'as r\'apido que CUDA cl\'asico ($\sim$2\% de diferencia) porque en este proyecto ambos usan kernels SIMT con algoritmos equivalentes.

\subsection{Comparativa GPU: RTX 4090 (servidor) vs RTX 3060 (local)}

\begin{table}[H]
\centering
\caption{Comparaci\'on de rendimiento GPU entre arquitecturas ($k=5$, $s=0.5$)}
\label{tab:gpu_comp}
\begin{tabular}{lccc}
\toprule
\textbf{Instancia} & \textbf{RTX 4090 (ms)} & \textbf{RTX 3060 (ms)} & \textbf{4090 m\'as r\'apida} \\
\midrule
Small    & 0.63 & 2.46 & 3.9$\times$ \\
Medium   & 4.24 & 8.72 & 2.1$\times$ \\
Large    & 14.55 & 25.05 & 1.7$\times$ \\
No-div   & 2.08 & 5.04 & 2.4$\times$ \\
\bottomrule
\end{tabular}
\end{table}

La RTX 4090 es entre 1.7$\times$ y 3.9$\times$ m\'as r\'apida que la RTX 3060. La diferencia es mayor en im\'agenes peque\~nas (3.9$\times$) donde la 4090 reduce el overhead de lanzamiento de kernel gracias a su mayor frecuencia y ancho de banda. En im\'agenes grandes, el speed-up se acerca al factor te\'orico dado por la diferencia en CUDA Cores (16,384 vs 3,584 $\approx$ 4.6$\times$), limitado por el ancho de banda de memoria.

\subsection{Efecto del Tama\~no del Kernel ($5\times5$ vs $9\times9$, servidor, medium)}

\begin{table}[H]
\centering
\caption{Efecto del kernel gaussiano en servidor (medium, $s=0.5$)}
\label{tab:kernel_effect}
\begin{tabular}{lccc}
\toprule
\textbf{Kernel} & \textbf{CPU (ms)} & \textbf{CUDA (ms)} & \textbf{Speed-up} \\
\midrule
$5 \times 5$ (25 ops/px)  & 317.3 & 4.24 & 74.9$\times$ \\
$9 \times 9$ (81 ops/px)  & 725.6 & 4.33 & 167.5$\times$ \\
\bottomrule
\end{tabular}
\end{table}

El kernel $9\times9$ implica 3.24$\times$ m\'as multiplicaciones que el $5\times5$. En CPU el slowdown es de 2.3$\times$, pero en GPU es solo 1.02$\times$: la GPU paraleliza el trabajo extra casi sin costo adicional. El speed-up pasa de 74.9$\times$ a 167.5$\times$.

\subsection{Convoluci\'on Directa vs Separable (servidor, medium, $k=5$)}

\begin{table}[H]
\centering
\caption{Comparaci\'on directa vs separable en servidor}
\label{tab:separable}
\begin{tabular}{lcccc}
\toprule
\textbf{M\'etodo} & \textbf{Ops/px} & \textbf{T$_{\text{kernel}}$ (ms)} & \textbf{T$_{\text{total}}$ (ms)} & \textbf{Throughput (MP/s)} \\
\midrule
Directa 2D   & 25 & 0.127 & 4.236 & 33,085 \\
Separable    & 10 & 0.095 & 4.310 & 44,082 \\
\bottomrule
\end{tabular}
\end{table}

La versi\'on separable reduce el tiempo de kernel del blur en un 25\% (de 0.127 a 0.095 ms) y aumenta el throughput en un 33\%. Sin embargo, el tiempo total es pr\'acticamente id\'entico (4.236 vs 4.310 ms) porque la transferencia H$\leftrightarrow$D domina el tiempo total en im\'agenes de tama\~no medio.

\subsection{Throughput por Versi\'on (servidor, $k=5$, $s=0.5$)}

\begin{table}[H]
\centering
\caption{Throughput promedio en megap\'ixeles por segundo (servidor)}
\label{tab:throughput}
\begin{tabular}{lccccc}
\toprule
\textbf{Versi\'on} & \textbf{Small} & \textbf{Medium} & \textbf{Large} & \textbf{No-div} & \textbf{Promedio} \\
\midrule
CPU Secuencial  & 13 & 13 & 13 & 13 & 13 \\
CUDA Cl\'asico  & 7,420 & 37,747 & 45,595 & 26,149 & 25,887 \\
CUDA Tile       & 8,434 & 35,761 & 42,645 & 25,272 & 28,088 \\
cuTile Python   & 1,148 & 8,789  & 6,040  & 5,775  & 5,600 \\
\bottomrule
\end{tabular}
\end{table}

El throughput de la GPU es $\sim$2,000$\times$ mayor que el de CPU. CUDA Tile supera ligeramente a CUDA cl\'asico en el promedio (28,088 vs 25,887 MP/s, +8.5\%). cuTile Python alcanza $\sim$5,600 MP/s, unas 5$\times$ menos que CUDA C++ debido al overhead de PyTorch en el lanzamiento de operaciones.

\subsection{Comparativa CPU: i9-14900KF (servidor) vs Ryzen 5 5600X (local)}

\begin{table}[H]
\centering
\caption{Comparaci\'on de rendimiento CPU ($k=5$, $s=0.5$)}
\begin{tabular}{lccc}
\toprule
\textbf{Instancia} & \textbf{i9-14900KF (ms)} & \textbf{Ryzen 5 5600X (ms)} & \textbf{Ratio} \\
\midrule
Small    & 19.85  & 33.46  & 1.69$\times$ \\
Medium   & 317.32 & 541.82 & 1.71$\times$ \\
Large    & 1333.27& 2138.76& 1.60$\times$ \\
No-div   & 119.60 & 209.35 & 1.75$\times$ \\
\bottomrule
\end{tabular}
\end{table}

El i9-14900KF del servidor es consistentemente $\sim$1.7$\times$ m\'as r\'apido que el Ryzen 5 5600X local en todas las instancias. Esto se debe al mayor IPC, frecuencias m\'as altas, y la ventaja de la memoria DDR5 de 125 GB frente a la DDR4 de 32 GB.

\subsection{Gr\'aficos Comparativos (Servidor RTX 4090)}

\begin{figure}[H]
\centering
\begin{tikzpicture}
\begin{axis}[
    ybar,
    bar width=12pt,
    width=0.92\textwidth,
    height=7cm,
    ylabel={Speed-up ($\times$)},
    symbolic x coords={Small, Medium, Large, No-div},
    xtick=data,
    nodes near coords,
    every node near coord/.append style={font=\tiny},
    ymin=0, ymax=110,
    legend style={at={(0.5,-0.18)}, anchor=north, legend columns=2},
    grid=major,
]
\addplot[blue!60, fill=blue!30] coordinates {
    (Small,31.7) (Medium,74.9) (Large,91.7) (No-div,57.4)
};
\addplot[red!60, fill=red!30] coordinates {
    (Small,32.5) (Medium,75.8) (Large,91.9) (No-div,57.4)
};
\legend{CUDA Cl\'asico, CUDA Tile}
\end{axis}
\end{tikzpicture}
\caption{Speed-up GPU vs CPU en servidor (RTX 4090, $k=5$, $s=0.5$). El speed-up escala con el tama\~no de imagen. CUDA Tile es $\sim$2\% m\'as r\'apido que CUDA cl\'asico.}
\label{fig:speedup}
\end{figure}

\begin{figure}[H]
\centering
\begin{tikzpicture}
\begin{axis}[
    ybar,
    bar width=18pt,
    width=0.92\textwidth,
    height=6cm,
    ylabel={Throughput (MP/s)},
    symbolic x coords={CPU, CUDA Clasico, CUDA Tile, cuTile Python},
    xtick=data,
    nodes near coords,
    nodes near coords align={vertical},
    ymin=0, ymax=35000,
    legend style={at={(0.5,-0.18)}, anchor=north, legend columns=1},
    grid=major,
    point meta=explicit symbolic,
]
\addplot[green!50!black, fill=green!30] coordinates {
    (CPU,13) (CUDA Clasico,25887) (CUDA Tile,28088) (cuTile Python,5600)
};
\end{axis}
\end{tikzpicture}
\caption{Throughput promedio por versi\'on en servidor ($k=5$, $s=0.5$). La GPU supera a la CPU por $\sim$2,000$\times$. cuTile Python es $\sim$5$\times$ m\'as lento que CUDA C++ por overhead de PyTorch.}
\label{fig:throughput}
\end{figure}

\subsection{Im\'agenes de Salida y Entregables}

Cada ejecuci\'on de la pipeline genera cuatro im\'agenes PNG por etapa de procesamiento, almacenadas en \texttt{results/<versi\'on>/<instancia>/}:

\begin{itemize}
  \item \texttt{<imagen>\_gray.png} --- salida de la conversi\'on RGB $\rightarrow$ luminancia.
  \item \texttt{<imagen>\_blur\_k<K>.png} --- desenfoque gaussiano con kernel $K \times K$.
  \item \texttt{<imagen>\_sobel\_k<K>.png} --- detecci\'on de bordes de Sobel.
  \item \texttt{<imagen>\_resize\_s<S>.png} --- redimensionamiento bilineal con factor de escala $S$.
\end{itemize}

\subsection{Muestra Visual de la Pipeline (servidor, medium, $k=5$, $s=0.5$)}

\begin{figure}[H]
\centering
\begin{subfigure}{0.48\textwidth}
  \centering
  \includegraphics[width=\textwidth]{ak-47_gray.png}
  \caption{Original $\rightarrow$ Grayscale}
\end{subfigure}
\hfill
\begin{subfigure}{0.48\textwidth}
  \centering
  \includegraphics[width=\textwidth]{ak-47_blur_k5.png}
  \caption{Gaussian Blur ($5 \times 5$)}
\end{subfigure}

\vspace{0.3cm}

\begin{subfigure}{0.48\textwidth}
  \centering
  \includegraphics[width=\textwidth]{ak-47_sobel_k5.png}
  \caption{Operador de Sobel}
\end{subfigure}
\hfill
\begin{subfigure}{0.48\textwidth}
  \centering
  \includegraphics[width=\textwidth]{ak-47_resize_s0.500000.png}
  \caption{Resize Bilineal ($0.5 \times$, $1024^2$)}
\end{subfigure}

\caption{Pipeline completa sobre \texttt{ak-47.png} ($2048 \times 2048$) ejecutada en el servidor (RTX 4090, CUDA cl\'asico). Cada etapa transforma progresivamente la imagen de entrada.}
\label{fig:pipeline_output}
\end{figure}

\subsection{Formato del CSV de Resultados}

El archivo \texttt{results/resultados.csv} contiene 23 columnas con los datos crudos de cada repetici\'on:

\begin{table}[H]
\centering
\caption{Columnas del CSV de resultados}
\small
\begin{tabular}{lp{0.55\textwidth}}
\toprule
\textbf{Columna} & \textbf{Significado} \\
\midrule
\texttt{Version} & CPU\_Secuencial, CUDA\_Clasico, CUDA\_Tile, cuTile\_Python \\
\texttt{Instancia} & small, medium, large, no-divisible \\
\texttt{Imagen} & nombre del archivo de entrada \\
\texttt{DimensionesOrig} & $W \times H$ originales \\
\texttt{KernelSize} & tama\~no del kernel gaussiano (5 \'o 9) \\
\texttt{Scale} & factor de resize (0.5 \'o 1.75) \\
\texttt{Bloque} & N/A (CPU), Dinamico\_API (CUDA), 16x16(SIMT\_en\_Tile), 16x16(cuTile) \\
\texttt{Repeticion} & 1 a 10 \\
\texttt{T\_\{Gray,Blur,Sobel,Resize\}\_kernel\_ms} & tiempo de kernel por etapa \\
\texttt{T\_\{...\}\_HtoD\_ms / DtoH\_ms} & tiempo de transferencia por etapa \\
\texttt{T\_Total\_ms} & suma de 4 etapas (incluye transferencias) \\
\texttt{Throughput\_MPps} & megap\'ixeles por segundo \\
\texttt{Herramienta} & chrono, CUDA\_Events, torch.cuda.Event \\
\bottomrule
\end{tabular}
\end{table}

El CSV contiene 800 filas de datos: 4 instancias $\times$ 2 kernels $\times$ 2 escalas $\times$ 10 repeticiones $\times$ 5 configuraciones (CPU, CUDA directo, CUDA separable, Tile, cuTile). La I/O de im\'agenes utiliza \texttt{stb\_image.h} / \texttt{stb\_image\_write.h} (single-header libraries de dominio p\'ublico) y Pillow para la versi\'on Python.

%=============================================================================
\section{An\'alisis T\'ecnico y Conclusiones}
%=============================================================================

\subsection{Escalabilidad y Ley de Amdahl}

El speed-up de CUDA sobre CPU escala con el tama\~no de imagen: desde 31.7$\times$ en $512^2$ hasta 91.7$\times$ en $4096^2$ en el servidor. Esto es esperable seg\'un la Ley de Amdahl: a mayor n\'umero de p\'ixeles, mayor es la fracci\'on paralelizable del trabajo total.

En la RTX 3060 local, el speed-up m\'aximo fue de 156.5$\times$ para large con $k=9$, mientras que en la RTX 4090 del servidor fue de 167.5$\times$ para la misma configuraci\'on. La diferencia se explica porque el speed-up de la 4090 est\'a limitado por el hecho de que la CPU del servidor (i9-14900KF) es $\sim$1.7$\times$ m\'as r\'apida que la CPU local (Ryzen 5 5600X), reduciendo el cociente.

\subsection{Convoluci\'on Separable: Ganancia Te\'orica vs Pr\'actica}

La versi\'on separable demuestra emp\'iricamente la reducci\'on de $O(k^2)$ a $O(2k)$. Para $k=5$, el tiempo de kernel del blur baja un 25\% (0.127 $\rightarrow$ 0.095 ms). Sin embargo, en la pr\'actica el beneficio es limitado porque:

\begin{enumerate}
  \item La transferencia H$\leftrightarrow$D domina el tiempo total en im\'agenes peque\~nas y medianas ($\sim$80\% del tiempo total en small).
  \item El buffer intermedio ($d\_temp$) agrega una asignaci\'on extra de memoria GPU.
  \item La sincronizaci\'on entre los dos pases introduce latencia adicional.
\end{enumerate}

Para kernels muy grandes ($k \geq 9$) y en im\'agenes grandes donde el c\'omputo domina sobre la transferencia, la versi\'on separable ofrecer\'ia una ventaja m\'as significativa, reduciendo 81 a 18 operaciones por p\'ixel (4.5$\times$ menos).

\subsection{CUDA Cl\'asico vs CUDA Tile}

En este proyecto, CUDA Tile es solo marginalmente m\'as r\'apido que CUDA cl\'asico ($\sim$2\% de diferencia, 28,088 vs 25,887 MP/s). Esto se debe a que ambos usan kernels SIMT con el mismo algoritmo. El flag \texttt{-enable-tile} compila correctamente y contiene un \texttt{\_\_tile\_global\_\_} real de demostraci\'on, pero la API actual de CUDA 13.3 no soporta stencils con acceso a vecinos, lo que impide explotar la ventaja te\'orica de Tile para Sobel/Gauss.

\subsection{Cuellos de Botella Identificados}

El an\'alisis con Nsight Systems 2026.1.3 sobre el binario CUDA cl\'asico (\texttt{medium}, $k=5$, $s=0.5$) revela la siguiente distribuci\'on del tiempo:

\begin{table}[H]
\centering
\caption{Distribuci\'on temporal de la pipeline en GPU (Nsight Systems, medium, $k=5$)}
\begin{tabular}{lcc}
\toprule
\textbf{Operaci\'on} & \textbf{Tiempo (ms)} & \textbf{\% del total} \\
\midrule
\texttt{cudaMalloc} (4 buffers) & 0.31 & 2.6\% \\
Transferencia H$\rightarrow$D (datos + kernel) & 2.48 & 21.3\% \\
Kernels (gray + blur + sobel + resize) & 1.50 & 12.9\% \\
Transferencia D$\rightarrow$H (resultados) & 3.47 & 29.8\% \\
\texttt{cudaFree} (4 buffers) & 0.05 & 0.4\% \\
Overhead de sincronizaci\'on y API & 3.81 & 33.0\% \\
\midrule
\textbf{Total} & \textbf{11.62} & \textbf{100\%} \\
\bottomrule
\end{tabular}
\end{table}

Las observaciones clave del perfil de Nsight Systems son:

\begin{enumerate}
  \item \textbf{Transferencia H$\leftrightarrow$D domina en im\'agenes peque\~nas:} en \texttt{small} ($512^2$), $\sim$80\% del tiempo total es transferencia por PCIe. En \texttt{large} ($4096^2$) la fracci\'on de c\'omputo aumenta al $\sim$40\%.
  \item \textbf{\texttt{cudaMalloc}/\texttt{cudaFree} por iteraci\'on:} cada wrapper de etapa reserva y libera memoria GPU dentro del bucle de medici\'on. Preasignar buffers una sola vez fuera del bucle reducir\'ia $\sim$0.3 ms de overhead fijo por iteraci\'on, un 3\% del tiempo total.
  \item \textbf{Convoluci\'on 2D directa:} aunque la versi\'on separable mitiga parcialmente el costo del blur ($0.127 \rightarrow 0.095$ ms), la pipeline a\'un se beneficiar\'ia de shared memory para que hilos vecinos reutilicen p\'ixeles y reduzcan accesos a memoria global.
\end{enumerate}

\textbf{Nota sobre Nsight Compute:} la herramienta \texttt{ncu} (Nsight Compute 2026.2.0) est\'a correctamente invocada en \texttt{profile.sh} con los flags \texttt{--set full --target-processes all}. Sin embargo, en ambos entornos de prueba el runtime de NCU requiere permisos de administrador para acceder a los GPU Performance Counters (\texttt{ERR\_NVGPUCTRPERM}). En Windows esto se resuelve ejecutando PowerShell como administrador; en Linux configurando \texttt{sudo ncu} o habilitando el registro de performance counters. Los flags y comandos est\'an documentados en \texttt{profile.sh} / \texttt{profile.bat} para ejecuci\'on en un entorno con los permisos adecuados.

\subsection{Determinismo y Reproducibilidad}

Los resultados son plenamente reproducibles entre ambas m\'aquinas: la misma pipeline, con los mismos par\'ametros e im\'agenes de entrada, produce las mismas im\'agenes de salida (verificado visual y num\'ericamente con MAE $\leq$ 0.08). Las diferencias de tiempo son exclusivamente atribuibles a las diferencias de hardware (n\'umero de CUDA cores, frecuencia de reloj, ancho de banda de memoria).

\subsection{Conclusi\'on Final}

\begin{enumerate}
  \item \textbf{Aceleraci\'on masiva:} la GPU reduce el tiempo de procesamiento de una imagen $4096 \times 4096$ de 21.8 minutos (CPU local) a 14.6 milisegundos (GPU servidor), un speed-up de 91.7$\times$ que demuestra que la pipeline de procesamiento de im\'agenes encaja naturalmente en el modelo SIMT de CUDA.
  \item \textbf{Convoluci\'on separable:} implementada y validada, reduce el tiempo de kernel del blur en 25\% adicional, demostrando que la optimizaci\'on algor\'itmica complementa la aceleraci\'on por hardware.
  \item \textbf{Escalabilidad cross-platform:} la misma base de c\'odigo compila y ejecuta en Windows y Linux, en GPUs desde RTX 3060 hasta RTX 4090, produciendo resultados id\'enticos. Los scripts \texttt{build.sh}, \texttt{run\_all.sh} y \texttt{profile.sh} garantizan reproducibilidad completa en Linux.
  \item \textbf{Tile C++:} el toolchain Tile est\'a habilitado y funcional, aunque limitado por la API actual de CUDA 13.3 para stencils. La mezcla de Tile + SIMT en la misma unidad de traducci\'on es una soluci\'on pr\'actica y oficialmente soportada.
\end{enumerate}

%=============================================================================
\section{Limitaciones y Trabajo Futuro}
%=============================================================================

\begin{itemize}
  \item \textbf{Imagen no-divisible:} la instancia utilizada es $1402 \times 1122$, no exactamente $1537 \times 1021$ como sugiere el enunciado. Ambas cumplen el requisito de no ser m\'ultiplos de 16, 32 ni 64, validando igualmente el manejo de dimensiones no alineadas a tiles.
  \item \textbf{API Tile C++:} la API \texttt{cuda\_tile.h} de CUDA 13.3 no soporta stencils con acceso a vecinos de forma pr\'actica. Cuando NVIDIA habilite esta funcionalidad, los 4 kernels de la pipeline podr\'an migrarse a \texttt{\_\_tile\_global\_\_} puro para una comparaci\'on m\'as estricta.
  \item \textbf{cuTile Python --- padding:} el kernel \texttt{rgb\_to\_gray\_kernel} de cuTile Python utiliza \texttt{PaddingMode.ZERO} en los tiles de borde, mientras que las dem\'as versiones usan clamp. La diferencia es insignificante para el c\'alculo de luminancia (el padding solo afecta p\'ixeles fuera de la imagen, que no se almacenan en la salida).
  \item \textbf{Convoluci\'on gaussiana separable:} se implement\'o solo para CUDA cl\'asico. Extenderla a CUDA Tile C++ y cuTile Python permitir\'ia una comparaci\'on completa del beneficio de $O(2k)$ vs $O(k^2)$ en todas las versiones.
  \item \textbf{Reutilizaci\'on de buffers GPU:} los wrappers CUDA reservan y liberan memoria en cada iteraci\'on. Preasignar buffers una vez y reutilizarlos a lo largo de las 10 repeticiones reducir\'ia el overhead de \texttt{cudaMalloc}/\texttt{cudaFree}, beneficiando especialmente a im\'agenes peque\~nas.
  \item \textbf{Shared memory:} los kernels actuales leen directamente de memoria global. Para Sobel y Gauss, cargar una regi\'on de p\'ixeles en shared memory por bloque reducir\'ia los accesos redundantes a memoria global (cada p\'ixel es le\'ido hasta $k^2$ veces por hilos vecinos).
\end{itemize}

\end{document}
